\FloatBarrier
\appendix
\section{Fitting the V12 mixture model}
\label{app:v12-fitting}

For a given Rubin band, let \(m\) be the S\'{e}rsic magnitude and let
\(x=\Delta m=m_{\mathrm{PSF}}-m_{\mathrm{S\acute{e}rsic}}\).
At each magnitude \(m\), the stellar profile \(U(x\mid m)\) and the
galaxy profile \(R(x\mid m)\) describe the distribution of \(x\) for
stars and galaxies, respectively. Each profile integrates to one. If
\(\omega_U(m)\) is the stellar fraction at magnitude \(m\), the
distribution of \(x\) for all objects is
\begin{equation}
 p(x\mid m)=\omega_U(m)\,U(x\mid m)+[1-\omega_U(m)]\,R(x\mid m).
 \label{eq:v12-mixture}
\end{equation}

We use HST labels to identify reference stars and galaxies
across the full COSMOS footprint. In Step~1, we fit the stellar profile
\(U(x\mid m)\) to Rubin measurements of the reference stars. In Step~2,
we fit the galaxy profile \(R(x\mid m)\) to Rubin measurements of the
reference galaxies. To apply the resulting \(r\)-band model to the
\(1^\circ\times1^\circ\) COSMOS field used in this paper, we keep both
profiles fixed at their full-COSMOS fits. In Step~3, we use the Rubin
measurements in this field to fit the stellar fraction \(\omega_U(m)\)
without using HST labels.

For each single-band fit, we use objects with a finite \(\Delta m\)
and a S\'{e}rsic magnitude \(18 \leq m < 26\) in the band being fitted.
This cut defines the sample used to fit that band's model.

For multi-band fit, we select objects with \(18 \leq r_{\mathrm{S\acute{e}rsic}} < 26\) from the original Rubin catalog.
We require finite, positive PSF and S\'{e}rsic fluxes in all four bands
\(g,r,i,z\).

\subsection{Step 1. Fit the stellar profile \texorpdfstring{\(U\)}{U}}

We use HST labels to identify the reference stars matched to
Rubin objects. For each reference star, the fit uses its Rubin
S\'{e}rsic magnitude \(m\) and its Rubin morphology measurement
\(x=\Delta m\) in the band being fitted.

The stellar profile combines a Johnson \(S_U\) component \(J\), which
describes the main concentration of stars, and a shifted lognormal
component \(T\), which describes the right tail:
\begin{equation}
 U(x\mid m)=[1-f(m)]\,J(x\mid m)+f(m)\,T(x\mid m),
 \label{eq:v12-stellar-profile}
\end{equation}
where \(f(m)\) is the fraction of the total area under \(U\) assigned
to \(T\) at magnitude \(m\). For example, \(f(m)=0.2\) assigns 20\%
of the stellar probability to \(T\) and 80\% to \(J\). We fit \(\operatorname{logit} f(m)\) jointly with the other stellar
profile parameters and convert it back to \(f(m)\). This keeps
\(f(m)\) between zero and one.

Let \(\phi(z)=(2\pi)^{-1/2}\exp(-z^2/2)\) be the standard normal
density. Its cumulative distribution function
\(\Phi(z)=\int_{-\infty}^{z}\phi(t)\,\mathrm{d}t\) gives the area
under the standard normal curve to the left of \(z\).
The cumulative distributions of \(J\) and \(T\) are
\begin{align}
 F_J(x\mid m) &=
 \Phi\!\left[\gamma+\delta\operatorname{asinh}
                 \!\left(\frac{x-\mu}{\sigma_c}\right)\right],\\
 F_T(x\mid m) &=
 \begin{cases}
 0, & x\leq\mu,\\[.15em]
 \displaystyle
 \Phi\!\left[\frac{1}{h}
             \ln\!\left(\frac{x-\mu}{\sigma_c\kappa}\right)\right],
                   & x>\mu.
 \end{cases}
 \label{eq:v12-stellar-cdf}
\end{align}
The two components share the location parameter \(\mu\). For \(J\),
\(\sigma_c\) sets the scale, and \(\gamma\) and \(\delta\) control
the shape. For \(T\), \(\sigma_c\kappa\) sets the scale of the
positive offset \(x-\mu\), and \(h\) sets its spread on a log scale.

We fit seven quantities as functions of magnitude.
\begin{equation}
 x_J,\quad \ln\sigma_c,\quad \ln(-\gamma),\quad
 \ln\delta,\quad\ln h,\quad\operatorname{logit}f,\quad
 \eta=\ln\!\left(\frac{x_T-\mu}{x_J-\mu}\right).
 \label{eq:v12-stellar-parameters}
\end{equation}
The position \(x_J\) specifies the peak of \(J\), while \(\sigma_c\),
\(\gamma\), and \(\delta\) control its scale and shape. For \(T\),
\(h\) controls the spread on a log scale, and \(\eta\) determines
where its peak lies relative to the peak of \(J\). The fraction \(f\)
controls how much of the total stellar probability goes into \(T\).
We fit the logit of the tail fraction,
\(\operatorname{logit}f=\ln[f/(1-f)]\), as a function of magnitude.
We then convert it back to \(f\), which gives the fraction of the
stellar profile assigned to \(T\) and stays between zero and one.

At each magnitude, we first determine the shared location \(\mu\).
The shape of \(J\) determines how far its peak lies to the right
of \(\mu\). We calculate this distance by solving for the positive
value of \(y\) in
\begin{equation}
 \delta y+\frac{\tanh y}{\delta}=-\gamma .
\end{equation}
The peak of \(J\) lies a distance \(\sigma_c\sinh y\) to the right
of \(\mu\). Since the fitted peak position is \(x_J\), we obtain
\(\mu\) by moving left from \(x_J\) by that distance.

Once we have \(\mu\), the fitted \(\eta\) tells us where the peak of
\(T\) should lie. The distance \(x_T-\mu\) is \(e^\eta\) times the
distance \(x_J-\mu\). For example, if \(\eta=\ln 2\), the peak of
\(T\) lies twice as far from \(\mu\) as the peak of \(J\).

Knowing \(x_T\) tells us where \(T\) should peak. To evaluate the
lognormal profile itself, we also need its scale
\(\sigma_c\kappa\). The peak position of a shifted lognormal depends
on both its scale and its log width \(h\). With \(h\) already
specified by the fit, we calculate \(\kappa\) so that the resulting
profile peaks at the position set by \(\eta\). The three derived
quantities are given by
\begin{equation}
 \mu=x_J-\sigma_c\sinh y,\qquad
 \ln\kappa=\ln(\sinh y)+h^2+\eta,\qquad
 x_T=\mu+\sigma_c\kappa e^{-h^2}.
 \label{eq:v12-stellar-derived}
\end{equation}

We find the peak of the full stellar profile by combining \(J\) and
\(T\) with the weights in equation~\eqref{eq:v12-stellar-profile} and
finding where the combined density is highest. Its position therefore
also depends on \(f\).

In the single-band fit, we evaluate \(U(x_i\mid m_i)\) at each
reference star's measured \(x_i\) and \(m_i\). We adjust the seven
magnitude-dependent functions together to maximize
\(\sum_i\ln U(x_i\mid m_i)\) over the reference stars, recalculating
\(\mu\), \(\kappa\), and \(x_T\) for each trial fit. The fitted coefficients of \(U(x\mid m)\) will be fixed in Step~3.

\subsection{Step 2. Fit the galaxy profile \texorpdfstring{\(R\)}{R}}

We fit \(R\) to the Rubin measurements of the matched reference
galaxies, using the same magnitude and measurement requirements as in
Step 1. The galaxy profile combines two skew-normal components,
\begin{equation}
 R(x\mid m)=[1-v(m)]\,R_1(x\mid m)+v(m)\,R_2(x\mid m),
 \label{eq:v12-galaxy-profile}
\end{equation}
where \(v(m)\) is the fraction within \(R\) assigned to \(R_2\).
Component \(R_j\) has location \(\ell_j\), positive scale \(s_j\), and
skew parameter \(\alpha_j\). With \(z_j=(x-\ell_j)/s_j\), its density is
\begin{equation}
 R_j(x\mid m)=\frac{2}{s_j}\phi(z_j)\Phi(\alpha_j z_j),
\end{equation}
where \(\phi\) and \(\Phi\) are the standard normal density and
cumulative distribution function defined in Step~1.

The seven fitted quantities are
\begin{equation}
 \ell_1,\quad\ln s_1,\quad\alpha_1,\quad
 \ell_2,\quad\ln s_2,\quad\alpha_2,\quad\operatorname{logit}v.
\end{equation}
We use an ordering constraint on the two component means to keep the
labels \(R_1\) and \(R_2\) consistent. Each mean is calculated from that
component's location, scale, and skew. For the single-band fit, we maximize \(\sum_i \ln R(x_i\mid m_i)\)
over the matched reference galaxies.
The two class profiles from Steps 1 and 2 then enter Step 3.

\enlargethispage{2\baselineskip}
\begin{samepage}
\paragraph{Magnitude dependence in Steps 1 and 2.}
Each of the seven fitted quantities in either profile follows a cubic
spline. For a quantity \(a\), the form is
\begin{equation}
 \begin{split}
 a(m)={}&c_0+c_1t+c_2t^2+c_3t^3
      +c_4(t+1/3)_+^3+c_5(t-1/3)_+^3,\\
 t={}&(m-22)/4,\qquad z_+=\max(0,z).
 \end{split}
 \label{eq:v12-spline}
\end{equation}
We fit the values of each spline at six equally spaced magnitudes from
18 to 26. These six values specify the same spline as the coefficients
\(c_0,\ldots,c_5\). Each class fit therefore estimates 42 coefficients
using the density at each object's measured \((m_i,x_i)\). The HST
reference fit starts from its class histograms. The Rubin fit tries two
data-derived starting points and keeps the converged, feasible solution
with the highest likelihood.
\par
\end{samepage}


